\documentclass[12pt,a4paper]{article}
\usepackage[utf8]{inputenc}
\usepackage[spanish]{babel}
\usepackage{amsmath,amssymb,amsfonts}
\usepackage{geometry}
\usepackage{hyperref}
\geometry{margin=2.5cm}

\title{\textbf{Constantes Matemáticas en Cálculo Diferencial}}
\author{Compendio de Constantes Fundamentales}
\date{\today}

\begin{document}

\maketitle

\section*{Introducción}
El cálculo diferencial estudia el cambio, las tasas de variación instantáneas y las pendientes de curvas mediante el concepto central de derivada. Las constantes en este campo desempeñan roles cruciales como bases para funciones trascendentes, crecimiento continuo e invariantes bajo la operación de diferenciación.

\section{Número de Euler / Base del Logaritmo Natural ($e$)}

\subsection*{1. Significado, Símbolo y Explicación}
Designado por la letra $e$ en honor a Leonhard Euler, es un número trascendental e irracional de inmensa importancia en análisis matemático. Se define formalmente mediante el límite:
\begin{equation*}
e = \lim_{n \to \infty} \left(1 + \frac{1}{n}\right)^n = \sum_{k=0}^{\infty} \frac{1}{k!}
\end{equation*}
En cálculo diferencial, su propiedad fundamental radica en que la función exponencial $f(x) = e^x$ es su propia derivada:
\begin{equation*}
\frac{d}{dx}\left(e^x\right) = e^x
\end{equation*}
Esto implica que en el punto $(0,1)$, la pendiente de la recta tangente a la curva $y = e^x$ es exactamente igual a $1$.

\subsection*{2. Valor Típico en Ingeniería y Ciencia}
Utilizado extensamente en ecuaciones de crecimiento poblacional, decaimiento radiactivo, circuitos RC/RL y procesamiento de señales:
\begin{equation*}
e \approx 2.7183 \quad \text{o} \quad 2.71828183
\end{equation*}

\subsection*{3. Valor Exacto a 15 Decimales}
\begin{equation*}
e \approx 2.718281828459045
\end{equation*}

\subsection*{4. Valor Exacto a 100 Decimales}
\begin{align*}
e \approx \; & 2.71828182845904523536028747135266249775724709369995957496696762772407... \\
& ...663035354759457138217852516642742747
\end{align*}
\textbf{Margen de error:} Constante matemática exacta (sin margen de error).

\vspace{0.8cm}
\hrule
\vspace{0.8cm}

\section{Constante de Euler-Mascheroni ($\gamma$)}

\subsection*{1. Significado, Símbolo y Explicación}
Representada con la letra griega gamma ($\gamma$), se define como el límite de la diferencia entre la serie armónica y el logaritmo natural:
\begin{equation*}
\gamma = \lim_{n \to \infty} \left( \sum_{k=1}^{n} \frac{1}{k} - \ln(n) \right)
\end{equation*}
Aparece de forma recurrente al derivar la función Gamma de Euler y en las derivadas de las funciones digamma ($\psi(x)$) y funciones especiales en cálculo avanzado.

\subsection*{2. Valor Típico en Ingeniería y Ciencia}
\begin{equation*}
\gamma \approx 0.5772 \quad \text{o} \quad 0.57721566
\end{equation*}

\subsection*{3. Valor Exacto a 15 Decimales}
\begin{equation*}
\gamma \approx 0.577215664901533
\end{equation*}

\subsection*{4. Valor Exacto a 100 Decimales}
\begin{align*}
\gamma \approx \; & 0.57721566490153286060651209008240243104215933593992359880576723488486... \\
& ...7726777664670936947063291746749514631
\end{align*}
\textbf{Margen de error:} Constante matemática pura exactísima (sin incertidumbre experimental).

\end{document}
